\chapter{Data entry: typed SQL, CRUD, forms, and PRG}

A classic web-application cycle is: read data from the database, accept typed input through an HTML form, validate it, authorize the object or operation, mutate state in a transaction, then redirect and read again with GET. This chapter treats that entire cycle as one unit. The important distinction is that \emph{valid input is not yet authorized input}.

\section{Typed SQL and CRUD}

\subsection*{Query}

\begin{lstlisting}[language=RWLang]
query fn articleBySlug(db: Db, slug: Slug)
    -> Result<Article, DbError> sql {
    SELECT id, slug, title, body
    FROM articles
    WHERE slug = :slug
}
\end{lstlisting}

The SQL text is static and compiler-owned. A user-provided value may enter only as a typed bind, such as \code{:slug}.

\subsection*{Mutation in a transaction}

\begin{lstlisting}[language=RWLang]
query fn articleById(db: Db, id: Int)
    -> Result<Article, DbError> sql {
    SELECT id, slug, title, body
    FROM articles
    WHERE id = :id
}

query fn updateArticle(
    tx: Transaction,
    id: Int,
    title: String,
    body: String
) -> Result<Void, DbError> mutates Article by id sql {
    UPDATE articles
    SET title = :title, body = :body
    WHERE id = :id
}
\end{lstlisting}

The action loads the object, establishes authorization, and only then uses the authorized object key for the mutation:

\begin{lstlisting}[language=RWLang]
action fn update(
    ctx: ActionContext,
    db: Db,
    id: Int,
    title: String,
    body: String
) -> Result<Redirect, PageError> {
    let article = articleById(db, id)?;
    authorize article authenticated;
    transaction db {
        updateArticle(tx, article.id, title, body)?;
    }
    return Ok(redirect(adminArticles()));
}
\end{lstlisting}

For an ordinary local transaction, success commits and a known database failure rolls back. A \code{Transaction} capability cannot escape the block. For retry-sensitive critical work, capture the explicit transaction outcome instead of assuming every failure proves rollback; Chapter~\ref{ch:db-deep-dive} explains \code{CommitUnknown}.

\subsection*{Portable database schema}

V1 aims for patterns that can remain portable across SQLite, PostgreSQL, and MariaDB. Canonical text may be used as the portable representation for \code{Date}, \code{DateTime}, \code{Uuid}, and \code{Decimal}.

\subsection*{Migration}

The server does not migrate automatically at startup. The recommended flow is:

\begin{lstlisting}
migrate verify
migrate apply
migrate verify
\end{lstlisting}

Keep migration credentials separate from normal application-runtime credentials.

\section{Forms, validation, and PRG}

\subsection*{Named form schema}

\begin{lstlisting}[language=RWLang]
form ArticleForm {
    title<String>
    body<String>
    published<Bool>
    validate title length 2 200 body length 1 200000
}
\end{lstlisting}

Route:

\begin{lstlisting}[language=RWLang]
route articleCreate POST "/admin/articles"
    form ArticleForm
    auth role Editor
    => create;
\end{lstlisting}

The handler parameter order and types should follow the form fields.

\subsection*{Invalid submission}

For a named form, a missing required field, a field-level typed decode failure, or a declared validation failure results in 422 \emph{Unprocessable Content}. The runtime refills field values in escaped form, attaches field errors, and inserts a CSRF token into the form. An unknown or duplicate field is instead a 400 because it violates the request schema.

This intentionally differs from inline form or JSON schemas: there a typed decode/validation failure is a request-contract violation and returns \code{400}. Use a named form when you want to report the error on the same user-facing form with field-level feedback.

\subsection*{Successful POST and PRG}

After a successful mutation, return a redirect so the browser loads the destination with GET. This Post/Redirect/Get pattern eliminates the common accidental resubmission problem.

\subsection*{A complete simple form-action pattern}

\begin{lstlisting}[language=RWLang]
action fn create(
    ctx: ActionContext,
    db: Db,
    title: String,
    body: String,
    published: Bool
) -> Result<Redirect, PageError> {
    let articleSlug = slug(title);
    transaction db {
        insertArticle(tx, articleSlug, title, body)?;
    }
    return Ok(redirect(adminArticles()));
}
\end{lstlisting}

For a named form, a missing checkbox-like \code{Bool} field is interpreted as \code{false}.

\subsection*{A practical decision order}

For a state-changing form, use this order during design and review:

\begin{enumerate}
  \item decode and validate the request at the route/form boundary;
  \item load the target object if the operation concerns an existing record;
  \item establish object/permission/tenant authorization;
  \item perform the mutation inside the appropriate transaction or critical-operation contract;
  \item redirect with a typed GET route helper;
  \item invalidate only the caches that the mutation actually makes stale.
\end{enumerate}

This ordering prevents a common mistake: treating a well-formed identifier or form as proof that the caller may modify the referenced object.
